\section{The mean and covariance jets in the principal comparison}
\label{sec:principal-comparison-jets}
We prove the extra smallness used in
\eqref{eq:principal-primitive-comparison}. The two means need not
agree away from the flat configuration. At the flat configuration
their values vanish and their first derivatives differ only by the
local inverse error. This gives one more small factor for their
difference than for the covariance difference.

\subsection{Uniform derivatives of the two local kernels}
We first record why the kernel bounds are uniform in the period.
All row estimates in this subsection use the maximum of weighted
row and column sums from Section~\ref{sec:local-covariance}.
For a fixed derivative order $n$, suppose a family of precisions
$A(b)$ has uniformly positive lower bound $r>0$, uniform
exponential row bound, and derivatives through order $n$ with
weighted row bounds at most $C(x)$, where $C$ has upper exponent
zero. The logarithmic fields $b$ range over a ball of radius
$v(x)$, with upper exponent $-1/2$. These hypotheses hold for
the unmasked local kernels under the fixed path and positive
precision bounds used in Section~\ref{sec:mean-action-jets}.

The inverse derivative formula, in multilinear notation, is
\begin{equation}\label{eq:comparison-inverse-derivative}
 D_I A^{-1}=
 \sum_{(I_1,\ldots,I_k)}(-1)^k
     A^{-1}(D_{I_1}A)A^{-1}\cdots(D_{I_k}A)A^{-1}.
\end{equation}
The sum is over ordered partitions of the labelled set $I$ into
nonempty blocks. Differentiate $AA^{-1}=I$ and induct on $|I|$
to prove this formula: the first differentiated block is the
nonempty subset falling on $A$, and the induction partitions the
remaining labels. Submultiplicativity and the inverse decay
estimate of Section~\ref{sec:local-covariance} bound each term
by a fixed power of $C(x)$ and of the inverse row bound.
The number of terms depends only on $n$. Thus all inverse jets
have upper exponent zero.

The same argument applies to every principal window, with a
small reduction of the exponential weight when the windows vary.
Indeed, at some fixed weight $\gamma>0$ the ordered-partition
formula bounds each inverse derivative on every window by $K_n$.
An assembled column is one such inverse column extended by zero;
thus each of its entries is at most
$K_n e^{-\gamma\operatorname{dist}(a_i,a_j)}$. At weight
$\gamma/2$, both row and column sums are bounded by
\[
 K_n m\sum_{q\ge0}(2q+1)^4 e^{-\gamma q/2},
\]
where $m$ bounds the number of coordinates at an anchor. This
bound does not require two columns to use the same window.
All windows and coordinate projections remain fixed during
 differentiation. Symmetrizing preserves the estimate.

On the reference side the coupling is the matrix of the
unmasked curvature and predictor derivatives, restricted by the
fixed keep test to the active coordinates. Denote it by $D(b)$;
its precision is exactly $D(b)^*D(b)$. Away from zero this
coupling differs from the completed-sector coupling: it uses the
actual derivatives of curved plaquettes and of the predictor,
whereas the completed sector transports their flat derivatives.
We bound the reference derivatives separately. A derivative of
order $j$ of a plaquette derivative is a derivative of order
$j+1$ of its finite product of link matrices; the product rule
bounds it by the same finite path-length moments, now through
order $j+1$. The predictor contribution is treated with its
fixed-ball derivative bound through that order. Purchase these
bounds through $n+1$ when requesting $n$ coupling derivatives.
All coefficients are bounded by the absolute path sums already
used in Section~\ref{sec:mean-action-jets}. Summing their
exponentially weighted rows and columns, and then applying the
fixed keep and active projections, gives an exponent-zero bound
for each reference coupling derivative. The product rule gives
the precision derivative bound, with a finite sum over subsets
of the derivative labels.
The flat identity proved below identifies $D(0)^*D(0)$ with the
active principal block of the completed sector precision; hence
Section~\ref{sec:sector-gram} gives its lower bound $r$.
Its perturbation on the logarithmic ball is bounded by $C(x)v(x)$,
so it remains at least $r/2$ eventually. Thus the inverse argument
applies to the reference matrix as well as the completed one.

The reference panels have a different geometric size from the
sector windows. For $B\ge120$ they contain every active coordinate
whose parent anchor is within $\lfloor B/60\rfloor$ of the column
anchor. Each endpoint of a bond lies within one coarse step of
its parent anchor. The triangle inequality therefore bounds the
separation of any two such endpoints by $\lfloor B/60\rfloor+2$;
the reference middle-panel criterion follows from
$20\lfloor B/60\rfloor+80\le B$. These are the two panel
inclusions required here. The inverse-error proof uses only
this inner radius, not an equality of the two window systems.
Accordingly the reference error costs
$e^{-c\lfloor B/60\rfloor}$ in place of $e^{-cB}$.
Both have every negative upper exponent.

For clarity, these derivatives are taken with the forest, tags,
read windows and sector indicators fixed. The unmasked matrix
entries are products of path derivatives and parallel transports.
Their row costs are finite sums of the absolute path prices
multiplied by fixed powers of path length. The fixed-ball
predictor bounds in Section~\ref{sec:mean-action-jets} make those
costs exponent zero. On the shrinking ball, the change from the
flat precision is at most $C(x)v(x)$ in operator norm, by the
fundamental theorem of calculus. Eventually this is less than
$r/2$, so the same inverse estimates apply throughout the ball.
This proves the required neighborhood of positivity without
changing any sector indicator.

Both affine sources are products of these kernel rows, their
coupling matrices and the unmasked residual vector. Their
derivatives therefore have exponent-zero bounds as well. The
residual and its derivatives have the fixed-ball bounds
\eqref{eq:action-fixed-ball}. A row acting on a bounded vector
uses its absolute row sum, not the total number of lattice
coordinates. The preceding product estimates are consequently
uniform over all finite periods and all allowed forests and tags.

\subsection{The exact comparison at the flat configuration}
Let $A$ denote the full flat sector precision, padded by the
identity on inactive coordinates, and let $U$ denote its active
coordinate set. At a flat field all parallel transports are
identities. Differentiating a plaquette gives its signed exterior
difference, and differentiating the predictor noise gives the
flat coarse response with the same sign in both constructions.
The normalization of the signed average is the same
$|\mathcal S|^{-1/2}$ in both matrices. Evaluating on each colour
coordinate vector therefore identifies the restricted sector
coupling with the reference coupling. Taking its Gram matrix gives
the exact flat identity
\begin{equation}\label{eq:comparison-flat-precision}
                             A_{UU}=A_{\rm ref}.
\end{equation}
The padding contributes zero on $U$. Rows discarded by the
reference keep test already vanish in this coupling, so their
removal does not change the product.

Since inactive rows and columns of $A$ are identity rows and
columns, $A$ is the direct sum of $A_{UU}$ and an identity matrix.
Hence
\[
                         (A^{-1})_{UU}=A_{UU}^{-1}.
\]
This conclusion uses the actual direct-sum property; it would
not follow merely by compressing an arbitrary inverse.
Let $P_6$ be the sector window inverse and $P_9$ the reference
window inverse at the flat point. The window-error calculation
of Section~\ref{sec:local-covariance}, applied to their respective
windows, gives
\begin{equation}\label{eq:comparison-flat-window-error}
              \|(P_6)_{UU}-P_9\|_0\le t(x),
\end{equation}
where $t$ has every negative upper exponent. Indeed each error
is bounded by its fixed inverse-decay constant times
$e^{-c\lfloor B_x/60\rfloor}$, and the cutoff radius $B_x$ dominates every fixed
multiple of $\log\beta_f(x)$. Taking the sum of the two errors
preserves this property. Symmetrizing
\eqref{eq:comparison-flat-window-error} gives the same bound
for the flat covariance difference.

Write the actual and reference means in logarithmic coordinates as
\[
              m_6(b)=-P_6(b)^*s_6(b),\qquad
              m_9(b)=-P_9(b)^*s_9(b).
\]
On the actual side the formula is restricted to $U$ after the
matrix-vector product. Each source is its coupling transpose
applied to the residual. The residuals vanish at $b=0$: all
curvature strokes are zero, the predictor is the identity, and
all coarse guides are linear in zero strokes. Thus both sources
and both means vanish there.

The source derivatives agree after restriction:
\begin{equation}\label{eq:comparison-source-first}
                         (Ds_6(0)h)|_U=Ds_9(0)h.
\end{equation}
To see this, write $s=M^*a$. At zero, $a(0)=0$, so its derivative
is $M(0)^*Da(0)h$ and the derivative of $M$ contributes zero.
The two residual formulas are the same unmasked curvature and
noise strokes, after restricting the kept rows. On omitted rows
the coupling is zero, as in
\eqref{eq:comparison-flat-precision}. This proves
\eqref{eq:comparison-source-first}. The actual source derivative
also vanishes outside $U$, because the corresponding coupling
columns are zero.

Differentiating the mean products now gives the exact identity
\begin{equation}\label{eq:comparison-mean-first}
 D(m_9-m_6)(0)h=
                  \bigl((P_6(0))_{UU}-P_9(0)\bigr)^*Ds_9(0)h.
\end{equation}
Derivatives of $P_a$ multiply the zero source and disappear.
By the row bound and the exponent-zero source derivative bound,
\[
          \|D(m_9-m_6)(0)\|\le t(x)C(x).
\]
This proves a first-derivative comparison, not merely an equality
of quadratic forms.

\subsection{Transferring the estimates along a local variation}
We give the elementary Taylor argument used for this transfer.
Let $f:E\to F$ be $C^n$, $n\ge2$, on a neighborhood of the
closed ball of radius $v\le1$, and suppose
\[
 f(0)=0,\qquad \|Df(0)\|\le Ct,
 \qquad \|D^kf(b)\|\le C\quad(2\le k\le n,\ \|b\|\le v).
\]
Let $u$ be a $C^n$ chart about the parameter base point with
$\|u(0)\|\le v$ and $\|D^ku(0)\|\le v^k$ for $1\le k\le n$.
Taylor's formula along the segment from $0$ to $b$ gives
\[
 \|f(b)\|\le Ct\,v+Cv^2/2,
                       \qquad \|Df(b)\|\le C(t+v).
\]
For a derivative of $f\circ u$ of order $k\ge1$, a partition
with $r$ blocks contributes $D^rf(u(0))$ applied to the
derivatives of $u$ of those block sizes. Their product costs
$v^k$. When $r=1$ this is at most $C(t+v)v$; when $r\ge2$
it is at most $Cv^2$, since $k\ge2$ and $v\le1$.
There are at most $n!2^n$ partitions for any $k\le n$.
Together with the value estimate this proves
\begin{equation}\label{eq:comparison-quadratic-jet-gain}
 \|f\circ u\|_{J^n_0}
                  \le n!2^n\{C(t+v)v+Cv^2\}.
\end{equation}
Similarly, if $g(0)$ has norm at most $t$ and all its
positive-order derivatives on the ball are bounded by $2C$,
the mean-value theorem and the same partition formula give
\begin{equation}\label{eq:comparison-linear-jet-gain}
                \|g\circ u\|_{J^n_0}\le n!2^n(t+2Cv).
\end{equation}

Apply \eqref{eq:comparison-quadratic-jet-gain} to
$f=m_9-m_6$, using \eqref{eq:comparison-mean-first}.
Apply \eqref{eq:comparison-linear-jet-gain} to each entry of
the covariance difference, using
\eqref{eq:comparison-flat-window-error}. The common logarithmic
chart from Sections~\ref{sec:haar-chart} and
\ref{sec:mean-action-jets} supplies the required $u$.
For a support anchored at $o$, its radius is at most $D_x$.
The chart's derivative scale can be chosen as
\[
 v(x)=C_n\,2\theta_f(x)
                    (1+\operatorname{reach}(x))^p(1+D_x)^p.
\]
This follows by applying the weighted direction bound at $o$
to every read site and using the uniform derivatives of the
logarithm on the fixed chart. An ordered product of $k$ local
variations contributes the $k$th power of that scale.
It has upper exponent $-1/2$ and is eventually at most one.

\subsection{Returning from the local logarithm to the physical field}
The preceding comparison is made with blank local sector data.
We must identify its readings with the physical ones; global
agreement of the fields or of their sector indicators is not
available. Let $Y$ be the allocated coarse support, disjoint
from the sector's seed set. For each coordinate read by the
insertion, with parent anchor $p$, its allocation includes
\[
 \operatorname{halo}_{R+4}\{p\}\subset Y,
                    \qquad R=45(1+B).
\]
The field is healthy on $Y$, and the period and chart bounds
ensure that this read region has the local gauge chart of
Section~\ref{sec:haar-chart}. Fix that gauge $g$ once, at the
base configuration. For every nearby configuration $W$ define
\[
 b(W)_e=\begin{cases}
 \log((g\cdot W)_e),&\operatorname{dist}(p,e_-)\le R,\\
 0,&\operatorname{dist}(p,e_-)>R.
 \end{cases}
\]
The finite set of logarithms stays in its common chart on a
neighborhood of the base point. Thus $b$ is smooth there, and
$\exp b(W)=g\cdot W$ on the entire read halo. No assertion is
made about plaquettes crossing the boundary of its zero extension.
The derivative estimates used above hold for precisely this
extension: the entries outside the halo have zero derivatives.

Here is why this local agreement suffices. A seed in the read
halo would belong to $Y$, contrary to seed separation. The local
indicator tests therefore agree with the blank tests there.
The forest, signed choices, widths and fixed path data are the
same. Every entry used by a sector window, and every source
entry used in its local mean, reads at most $17B$ from the
column anchor. This fits inside $R$. In the reference construction
the middle and end panels and their field queries are contained
in the same allocated halo, with the extra four sites allowing
for the bond endpoints. For the reference construction this containment can be checked
numerically. The panel radius is at most $\lfloor B/20\rfloor$.
If $E$ denotes the extent of a single path query, the coupling
read radius is $E+10B+4$, and the covariance read radius is this
quantity plus $2\lfloor B/20\rfloor$. The mean adds the residual
read radius, the coupling range, and one more panel radius;
write the resulting radius as $S$. The fixed range inequalities
give $S+E\le36B$. Even an allowed displacement of the query
anchor by $7(1+B)$ therefore costs at most
$7(1+B)+36B\le45(1+B)=R$. These inequalities cover the
whole local product, including source entries in a mean row.
Consequently their restricted coupling
and precision entries, and their source entries, agree after
the local gauge change. Equal principal matrices have equal
inverses; their restricted matrix products also agree. This
proves equality of the mean readings and window covariance
blocks. It uses the local tests only, so sector indicators outside
the halo can be arbitrary.

More explicitly, write $O_j$ for the colour adjoint of the fixed
gauge at the parent of coordinate $j$. For a mean read at $k$,
and a covariance block between $k$ and $l$, the preceding
identification gives
\[
 O_k m_a^{\rm phys}(W)_k=m_a^{\rm loc}(b(W))_k,
 \qquad
 C_a^{\rm phys}(W)_{kl}
   =O_k^* C_a^{\rm loc}(b(W))_{kl}O_l,\qquad a=6,9.
\]
For the sector covariance its locality initially identifies the
column at $k$; symmetry gives the displayed row identity. The
reference identity reads that row directly. Thus the same
chart at $k$ compares any requested block, including $l$ outside
the insertion's support. The mean identity on the reference
side holds for every fixed fine field and auxiliary parameter:
its local affine source is the same one just identified.
All these equalities hold on a neighborhood, so their derivatives
through the required order agree as well.

A colour adjoint is an isometry, so returning from the local
mean does not change its norm. Each entry of $O_j$ has magnitude
at most one, so the covariance entry bound acquires at most
$(\dim\mathfrak g)^2$. The gauge is fixed while differentiating.
Consequently the actual unscaled mean difference
has upper exponent $-1$, whereas the covariance difference has
upper exponent $-1/2$, exactly as in
\eqref{eq:principal-primitive-comparison}.

The argument uses the same local histories, flags and reference
windows on both sides. It does not identify the two precisions
on nonflat fields. The hypotheses inherited from the earlier
kernel and chart sections are the common path bounds, their
local smooth formulas, and the positive precision estimates;
these remain part of the selected construction throughout.
